\begin{statementbox}
\textbf{C014\quad 椭圆焦点与切线综合性质（修订版）}

设椭圆
\[
C:\ \frac{x^2}{a^2}+\frac{y^2}{b^2}=1,
\qquad a>b>0,\quad c^2=a^2-b^2,
\]
焦点为 $F_1(-c,0),F_2(c,0)$，中心为 $O(0,0)$。
点 $P(x_0,y_0)\in C$，$l$ 为 $P$ 处切线：
\[
\boxed{\ l:\ \frac{x_0x}{a^2}+\frac{y_0y}{b^2}=1\ }.
\]
$F_i$ 到 $l$ 的垂足记为 $H_i$，距离 $d_i=F_iH_i$（$i=1,2$）；
四边形 $F_1H_1H_2F_2$ 的面积记为 $S$。

\begin{enumerate}
\item \textbf{焦点到切线的距离积（全椭圆成立）}：
\[
\boxed{d_1d_2=b^2}.
\]
尤其在长轴右顶点 $P(a,0)$，有 $d_1=a+c,d_2=a-c$，
故 $(a+c)(a-c)=b^2$，是最直观的记忆锚点。

\item \textbf{垂足共圆（全椭圆成立）}：两个垂足均在长半轴为半径的\textbf{辅助圆（主圆）}上，
\[
\boxed{OH_1=OH_2=a},\qquad
\boxed{x_{H_i}^2+y_{H_i}^2=a^2}.
\]
注意辅助圆的半径为 $a$，不是 $b$，也不是焦距 $c$。

\item \textbf{光学反射（非长轴顶点时按通常角平分线理解）}：
当 $y_0\ne0$ 时，$P$ 与两焦点组成非退化三角形；
切线 $l$ 平分 $\angle F_1PF_2$ 的\textbf{外角}，
法线（过 $P$ 垂直 $l$ 的直线）平分其\textbf{内角}。
因此由 $F_1$ 射向 $P$ 的光，沿椭圆切线反射后通过 $F_2$。
在 $P=(\pm a,0)$ 处按反射直线的极限位置理解，不把退化角当作普通三角形角平分线。

\item \textbf{角平分线的 $x$ 轴截距积（有条件）}：
设 $m$ 为内角平分线（法线），$l$ 为外角平分线（切线）。
若 $x_0y_0\ne0$，记 $M=m\cap Ox$、$N=l\cap Ox$，则
\[
\boxed{x_M=\frac{c^2x_0}{a^2}},\qquad
\boxed{x_N=\frac{a^2}{x_0}},\qquad
\boxed{x_Mx_N=c^2}.
\]
若 $x_0=0$，切线平行 $x$ 轴，$N$ 不存在；若 $y_0=0$，焦点三角形退化，不能使用角平分线表述。

\item \textbf{四边形面积最大值（必须分类）}：
\[
\boxed{
S_{\max}=\begin{cases}
2bc,&c\le b,\\
a^2,&c\ge b.
\end{cases}}
\]
两分支在 $c=b$ 时一致。特别地，原先不加条件写成
``$S_{\max}=a^2$'' \textbf{不正确}。
更精确地说，$c<b$ 时仅在上下顶点（水平切线）取到 $2bc$；
$c=b$ 时在上下顶点取到 $a^2$；$c>b$ 时可在四个对称切点取到 $a^2$。
当切线为长轴端点处的竖直切线时，$H_1=H_2=P$，面积为 $0$。
\end{enumerate}

\textbf{核心推导链}：椭圆定义 $PF_1+PF_2=2a$
$\Longrightarrow$ 反射性质
$\Longrightarrow$ 中位线定理得 $OH_i=a$
$\Longrightarrow$ 主圆的点幂定理得 $d_1d_2=b^2$。
\end{statementbox}
